\documentclass[manuscript,review,anonymous,natbib=false]{acmart}

\usepackage[style=acmnumeric]{biblatex}
\usepackage{CJKutf8}
\addbibresource{references.bib}

\AtBeginDocument{%
  \providecommand\BibTeX{{%
    Bib\TeX}}}

\setcopyright{acmlicensed}
\copyrightyear{2018}
\acmYear{2018}
\acmDOI{XXXXXXX.XXXXXXX}

\acmConference[Conference acronym 'XX]{Make sure to enter the correct
  conference title from your rights confirmation email}{June 03--05,
  2018}{Woodstock, NY}

\acmISBN{978-1-4503-XXXX-X/2018/06}

\newcommand{\fz}{FIELD / Zh\`{e}n}

\title{\fz{}: Designing an AI-Driven Generative Soundtrack for British / Taiwanese Contemporary Dance}

\author{Nick Rothwell}

\bibliography{references.bib}

\begin{document}
\begin{CJK*}{UTF8}{gbsn}

  \maketitle

\begin{quote}
\fz{} is a contemporary dance work for stage blending English and Taiwanese folk dance, music, and storytelling with AI, projection, and motion tracking to reimagine myth and ritual as joyful, embodied visions of our shared future.
\end{quote}

\section{Introduction}

\fz{}, a folk/contemporary dance project co-devised by Pell Ensemble\cite{noauthor_rebecca_nodate} in the UK and Anarchy Dance Theatre\cite{noauthor_anarchy_nodate} in Taiwan, combines a conventional choreographic making process with motion capture, machine learning (AI), generative digital visuals and a generative sound score. The AI system generates motion capture data in real time from camera feeds of the action on stage, and also analyses the movement to determine higher-level gestural components and emotional intent. This poster is concerned with the design and implementation of a software system which generates an interactive sound score mediated by a continuous data feed from the AI.

\section{Motivation}

The use of AI is central to the creative focus and narrative of the work. The work's theme incorporates a machine/human relationship being explored and shared live with an audience, and this is manifested through the visuals and the sound score. The piece also explores and incorporates aspects of English folk dance and Taiwanese temple dance, both of these forms traditionally involving live musicians, and this aspect of the performance can be recontextualised as a dynamic, generative score created with electronic/software instruments and mediated by an AI data feed. (Electronic scores in the guise of folk collaboration are not new---for example the Blue Ghost Flamenco project by Luz \& Mannion\cite{noauthor_blue_nodate}.)

The desire to immerse the audience in an evolving, responsive soundscape means that there should be no pausing of sound for resets or preset changes between scenes; the soundscape should be continuously evolving, whilst able to deliver the desired sonic score or environment required for each of the multiple scenes in the dance work. The ability to crossfade between musical sections could be considered necessary, but certainly not sufficient for a continuously engaging sonic experience. In addition, at different points in the performance the score is required to respond to different parts of the AI data feed, or the same parts with different degrees of effect, and again this response pattern needs to change gradually when moving from one scene to the next, amplifying some parts of the AI data feed whilst attenuating or even discarding others.

Accordingly, the configuration and behaviour of the sound score responds in time at two levels of scale: short-term, it responds to realtime AI data being generated by the performers on-stage, whilst longer term it gradually transforms between scenes and preset states. (There is no technical reason why preset changes could not be close to instantaneous---think of this preset management process as DAW automation\cite{shimazu_what_2023}---but for the purposes of the project they are most likely to be relatively slow.)

\section{Instrumentation}

\end{CJK*}

\printbibliography

\end{document}
